\chapter{Synchronisation temporelle}\label{ch:nw:time-synchronisation}

Le droit européen oblige une société qui pratique le trading algorithmique à haute fréquence à maintenir les horloges qui horodatent ses
ordres à moins de 100~microsecondes de l'UTC, et à horodater à la microseconde près. La spécification du \omterm{def:nw:time-synchronisation:protocols}{Network Time Protocol} part
du principe qu'une horloge informatique ordinaire peut dériver en fréquence de 15~parties par million~: quinze microsecondes gagnées
ou perdues à chaque seconde. Une telle horloge, réglée puis laissée à elle-même, sort de la règle en sept secondes. Rester dans
cette limite suppose une boucle de régulation qui tourne toute la journée, et en prouver le respect suppose un registre qu'il faut conserver.

Le livre a partout utilisé le temps~: les horodatages d'exécution et de réception du livre~1, les horodatages matériels du
livre~13, les captures du chapitre~5 à venir. Ce chapitre explique d'où vient le temps d'une société de trading, comment il est
transmis à chaque serveur à travers un réseau qui le retarde de quantités variables, et ce que les régulateurs demandent à une société de démontrer.

\section{Horloges, oscillateurs et dérive}

Une horloge compte les cycles d'un oscillateur. Si l'oscillateur va légèrement trop vite, l'horloge prend de l'avance~; si sa
fréquence elle-même varie (température, vieillissement), l'avance varie aussi.

\begin{definition}[Décalage d'horloge, écart de fréquence relatif]\label{def:nw:time-synchronisation:offset}
Le \emph{décalage d'horloge}\index{décalage d'horloge} $\vartheta(t)$ d'une horloge est l'écart entre l'heure qu'elle affiche et
l'heure de référence au même instant. Son \emph{écart de fréquence relatif}\index{ecart de frequence relatif@écart de fréquence relatif} $y_{\mathrm f}(t) =
d\vartheta/dt$ est la vitesse à laquelle le décalage croît, un nombre sans dimension exprimé en parties par million (ppm) ou par
milliard (ppb)~: une horloge avec $y_{\mathrm f} = 1$~ppm gagne une microseconde à chaque seconde.
\end{definition}

\begin{definition}[Écart-type d'Allan]\label{def:nw:time-synchronisation:adev}
L'\emph{écart-type d'Allan}\index{ecart-type d'Allan@écart-type d'Allan} $\sigma_y(\tau)$ d'une horloge est la racine de la moyenne quadratique de la
variation de sa fréquence relative moyenne entre deux intervalles successifs de longueur $\tau$, divisée par $\sqrt2$~:
$\sigma_y^2(\tau) = \tfrac12 \E\bigl[(\bar y_{k+1} - \bar y_k)^2\bigr]$. À partir d'échantillons de phase $x_i$ espacés de $\tau$, il
vaut $\sigma_y^2(\tau) = \E\bigl[(x_{i+2} - 2x_{i+1} + x_i)^2\bigr] / (2\tau^2)$.
\end{definition}

L'\omterm{def:nw:time-synchronisation:adev}{écart-type d'Allan} répond à la question que se pose réellement une société~: sur un intervalle de $\tau$, de combien peut-on se
fier à la vitesse de l'horloge~? Contrairement à l'écart-type ordinaire, il converge pour les marches aléatoires que subissent les
oscillateurs réels, et la pente de sa courbe en fonction de $\tau$, tracée en échelle log-log, identifie le bruit~: $-1$ pour un bruit de phase blanc (une horloge dont les lectures sont bruitées mais dont la vitesse est juste), $+\tfrac12$ pour une marche
aléatoire de fréquence (une vitesse qui dérive). La \Cref{fig:nw:time-synchronisation:adev} montre les deux.

\begin{definition}[Maintien (holdover)]\label{def:nw:time-synchronisation:holdover}
Le \emph{maintien (holdover)}\index{maintien (holdover)} désigne le fonctionnement d'une horloge qui a perdu sa référence et continue de compter le temps sur son propre oscillateur, à partir
de la dernière correction de fréquence qu'elle a reçue~; la durée du maintien est le temps pendant lequel elle reste dans une tolérance.
\end{definition}

\begin{proposition}[Durée du maintien]\label{prop:nw:time-synchronisation:holdover}
Une horloge qui entre en maintien avec un décalage $\vartheta_0$, une erreur de fréquence résiduelle $y_0$ et un vieillissement constant $D$ (variation de fréquence
par unité de temps) a pour décalage $\vartheta(t) = \vartheta_0 + y_0 t + \tfrac12 D t^2$. Elle reste dans une limite $L > |\vartheta_0|$ jusqu'à la
première racine positive de $|\vartheta(t)| = L$~; sans vieillissement, $t = (L - \vartheta_0)/y_0$ pour $y_0 > 0$.
\end{proposition}

\begin{proof}
Intégrer $d\vartheta/dt = y_0 + Dt$ à partir de $\vartheta(0) = \vartheta_0$.
\end{proof}

\section{Transfert de temps bidirectionnel~: NTP et PTP}

Un serveur apprend l'heure d'un maître en échangeant des messages avec lui. La difficulté tient à ce que l'heure d'arrivée d'un
message mélange le décalage des horloges et le délai du réseau~; deux messages en sens opposés permettent de les séparer, sous une hypothèse.

\begin{definition}[Network Time Protocol, Precision Time Protocol]\label{def:nw:time-synchronisation:protocols}
Le \emph{Network Time Protocol}\index{Network Time Protocol} (NTP) synchronise un client sur des serveurs par des échanges
requête-réponse sur UDP, typiquement toutes les quelques secondes à quelques minutes, avec des horodatages pris en logiciel. Le \emph{Precision Time
Protocol}\index{Precision Time Protocol} (PTP, IEEE~1588) synchronise les horloges d'un réseau sur un maître par des échanges fréquents, plusieurs par
seconde, conçu pour des horodatages pris dans le matériel réseau et pour des commutateurs qui participent au protocole.
\end{definition}

\begin{omfigure}
\begin{tikzpicture}[font=\footnotesize,>=Stealth]
\draw[thick] (0,0) -- (0,-4.4) node[below] {maître};
\draw[thick] (5,0) -- (5,-4.4) node[below] {esclave};
\draw[->,omDef,thick] (0,-0.5) node[left]{$t_1$} -- node[above,sloped]{Sync (et Follow\_Up avec $t_1$)} (5,-1.3) node[right]{$t_2$};
\draw[->,omThm,thick] (5,-2.4) node[right]{$t_3$} -- node[above,sloped]{Delay\_Req} (0,-3.3) node[left]{$t_4$};
\draw[->,gray] (0,-3.5) -- node[below,sloped,font=\scriptsize]{Delay\_Resp (avec $t_4$)} (5,-4.2);
\draw[<->,gray] (5.9,-0.5) -- node[right,font=\scriptsize]{$d_{\mathrm{ms}}$ plus décalage} (5.9,-1.3);
\draw[<->,gray] (-0.9,-2.4) -- node[left,font=\scriptsize]{$d_{\mathrm{sm}}$ moins décalage} (-0.9,-3.3);
\end{tikzpicture}

\omcaption{Échange requête-délai de PTP. $t_1$ et $t_4$ sont lus sur l'horloge du maître, $t_2$ et $t_3$ sur celle de l'esclave. La
première différence est le délai de propagation du maître vers l'esclave plus le décalage de l'esclave~; la seconde est le délai du
retour moins le décalage. Le temps s'écoule vers le bas.}\label{fig:nw:time-synchronisation:exchange}
\end{omfigure}

\begin{proposition}[Décalage bidirectionnel et son angle mort]\label{prop:nw:time-synchronisation:offset}
Soit le maître qui émet à $t_1$ (son horloge), l'esclave qui reçoit à $t_2$ et répond à $t_3$ (son horloge), le maître qui reçoit à
$t_4$, et soient $d_{\mathrm{ms}}$ et $d_{\mathrm{sm}}$ les délais à sens unique. L'estimation
\[
  \hat\vartheta = \tfrac12\bigl[(t_2 - t_1) - (t_4 - t_3)\bigr]
  \quad\text{vaut}\quad \vartheta + \tfrac12\,(d_{\mathrm{ms}} - d_{\mathrm{sm}}),
\]
et l'aller-retour $(t_4 - t_1) - (t_3 - t_2) = d_{\mathrm{ms}} + d_{\mathrm{sm}}$ ne dépend pas de $\vartheta$. La méthode
bidirectionnelle est exacte lorsque les deux sens prennent le même temps, et aucun échange de messages ne permet de mesurer leur écart.
\end{proposition}

\begin{proof}
$t_2 = t_1 + d_{\mathrm{ms}} + \vartheta$ et $t_4 = t_3 - \vartheta + d_{\mathrm{sm}}$~; substituer. Tout triplet $(\vartheta, d_{\mathrm{ms}},
d_{\mathrm{sm}})$ et $(\vartheta + a, d_{\mathrm{ms}} - a, d_{\mathrm{sm}} + a)$ produisent les mêmes quatre horodatages.
\end{proof}

\begin{definition}[Asymétrie de chemin, asservissement d'horloge]\label{def:nw:time-synchronisation:asym}
L'\emph{asymétrie de chemin}\index{asymétrie de chemin} d'un chemin de synchronisation est la différence $d_{\mathrm{ms}} - d_{\mathrm{sm}}$
entre ses deux délais à sens unique~; la moitié de cette différence apparaît, sans être détectée, dans chaque estimation de décalage bidirectionnel. Un \emph{asservissement d'horloge (servo)}\index{asservissement d'horloge (servo)} est la
boucle de régulation qui transforme les estimations successives de décalage en corrections de la phase et de la fréquence de l'esclave.
\end{definition}

L'asymétrie est fixée par les câbles (une fibre plus longue dans un sens), par les émetteurs-récepteurs et par les commutateurs qui
traitent différemment les deux sens~; la mise en file d'attente dans chaque sens ajoute une part variable, que le moyennage réduit et
que l'asservissement filtre. Une société trouve la part fixe par calibration contre une deuxième référence indépendante, puis la soustrait par configuration~: aucun protocole ne peut la trouver.

\omcode{code/firm/clocksync/firm_clocksync.py}{58}{68}{Un échange dans \texttt{firm.clocksync}~: mise en file d'attente dans les deux sens, bruit des horodatages
et correction des horloges transparentes lorsqu'il y en a.}\label{lst:nw:time-synchronisation:exchange}

\section{Horloges grand maître (grandmaster), références de temps GNSS et maintien}

\begin{definition}[Horloge grand maître, référence de temps GNSS]\label{def:nw:time-synchronisation:gm}
Une \emph{horloge grand maître}\index{horloge grand maître} (grandmaster clock) est l'horloge au sommet d'un réseau PTP, maîtresse de toutes les autres. Une \emph{référence de
temps GNSS}\index{référence de temps GNSS} est un récepteur d'un système mondial de navigation par satellite (GPS, Galileo et d'autres) utilisé comme
source d'UTC~: le signal de chaque satellite porte l'heure du système et son écart à l'UTC, tel que le maintient un laboratoire national.
\end{definition}

L'\omterm{def:nw:time-synchronisation:gm}{horloge grand maître} d'une société de trading est en général un boîtier dédié dans chaque cage~: une antenne satellite sur le toit du centre de données, un câble
jusqu'au récepteur, et un oscillateur que le récepteur discipline et qui assure le \omterm{def:nw:time-synchronisation:holdover}{maintien (holdover)} de l'heure lorsque le signal est perdu
(brouillage, antenne défaillante, câble coupé). La réglementation européenne accepte l'UTC fourni par un système satellitaire «~à condition que tout écart à l'UTC
soit pris en compte et retiré~». La norme de performance propre au GPS fixe la précision de l'écart à l'UTC qu'il diffuse~: 30 nanosecondes,
95~\% du temps. Le reste du budget d'erreur revient à la société~: le délai du câble d'antenne (une longueur fixe et mesurable), le récepteur, et
chaque saut entre l'\omterm{def:nw:time-synchronisation:gm}{horloge grand maître} et l'horodatage.

\begin{omfigure}
\begin{tikzpicture}[font=\footnotesize,b/.style={draw,rounded corners=2pt,minimum height=0.65cm,align=center,inner sep=3pt},>=Stealth]
\node[b,fill=gray!10] (sat) at (0,0) {satellites\\ (écart à l'UTC\\ dans le signal)};
\node[b,fill=omDef!15] (gm) at (3.9,0) {horloge grand maître\\ récepteur et\\ oscillateur};
\node[b,fill=omMeth!18] (bc) at (7.4,0) {horloge frontière\\ (commutateur)};
\node[b,fill=omMeth!10] (tc) at (7.4,-1.9) {horloge transparente\\ (commutateur)};
\node[b] (h1) at (10.8,0.6) {serveur~: horloge de la carte};
\node[b] (h2) at (10.8,-0.6) {serveur~: horloge de la carte};
\node[b] (h3) at (10.8,-1.9) {boîtier de capture};
\draw[->,dashed] (sat) -- node[above,font=\scriptsize,align=center]{antenne\\ câble} (gm);
\draw[->,thick] (gm) -- node[above,font=\scriptsize]{PTP} (bc);
\draw[->,thick] (bc.east) -- (h1.west);
\draw[->,thick] (bc.east) -- (h2.west);
\draw[->,thick] (bc.south) -- (tc.north);
\draw[->,thick] (tc) -- (h3);
\node[font=\scriptsize,align=left,anchor=north west] at (-1.0,-1.2) {erreurs le long de la chaîne~: l'écart\\ à l'UTC diffusé, le câble d'antenne,
  le récepteur,\\ chaque horloge frontière, l'asymétrie\\ résiduelle, l'horloge de la carte,\\ le passage à l'horloge de l'hôte};
\end{tikzpicture}

\omcaption{Chaîne de temps d'une cage. L'\omterm{def:nw:time-synchronisation:gm}{horloge grand maître} prend l'UTC auprès des satellites et sert le PTP~; une \omterm{def:nw:time-synchronisation:bc}{horloge frontière} est esclave d'un côté
et maîtresse de l'autre~; une \omterm{def:nw:time-synchronisation:bc}{horloge transparente} relaie les messages PTP et les corrige du temps passé à l'intérieur d'elle. Chaque maillon
de la chaîne ajoute au budget d'erreur des horodatages situés à son extrémité.}\label{fig:nw:time-synchronisation:chain}
\end{omfigure}

\section{Horloges frontière et horloges transparentes~: suppression de l'erreur du réseau}

\begin{definition}[Horloge frontière, horloge transparente]\label{def:nw:time-synchronisation:bc}
Une \emph{horloge frontière}\index{horloge frontière} (boundary clock) est un équipement réseau doté de sa propre horloge, esclave du PTP sur un port et maîtresse sur
ses autres ports~: elle termine un échange et en démarre de nouveaux, de sorte qu'aucun message PTP ne la traverse. Une \emph{horloge transparente}\index{horloge transparente}
(transparent clock) relaie les messages PTP et ajoute à un champ de correction de chacun le temps passé par le message à l'intérieur de l'équipement (son temps
de séjour), que les extrémités soustraient de leurs horodatages.
\end{definition}

La mise en file d'attente dans un commutateur est le plus grand délai variable d'un chemin de synchronisation, et elle n'est pas symétrique~: le message Sync peut attendre derrière
une rafale de données de marché alors que le Delay\_Req, qui circule dans l'autre sens, n'attend pas. Une \omterm{def:nw:time-synchronisation:bc}{horloge transparente} mesure cette attente et la signale~; une
\omterm{def:nw:time-synchronisation:bc}{horloge frontière} supprime le problème en ne laissant jamais un message faire la file à travers le commutateur. Les deux exigent que le matériel propre au commutateur sache horodater, et
les deux laissent intacte l'asymétrie fixe des câbles et des émetteurs-récepteurs.

\begin{omfigure}
\begin{tikzpicture}
\begin{axis}[width=11.5cm,height=5.6cm,ymode=log,xmin=0,xmax=600,ymin=3,ymax=200000,
  xlabel={secondes après le début},ylabel={plus grand $|\vartheta|$ sur chaque intervalle de 5~s (ns, log)},
  ytick={10,100,1000,10000,100000},yticklabels={10,100,1\,000,10\,000,100\,000},
  tick label style={font=\small},label style={font=\small},grid=major,grid style={gray!25},
  legend columns=2,legend style={font=\footnotesize,at={(0.5,-0.3)},anchor=north,draw=none,fill=none,
  /tikz/every even column/.append style={column sep=8pt}},table/col sep=comma]
\addplot[omThm,thick] table[x=t_s,y=sw]{figdata/networks/04-time-synchronisation/offset.csv};
\addplot[omExo,thick] table[x=t_s,y=hw]{figdata/networks/04-time-synchronisation/offset.csv};
\addplot[omMeth,thick] table[x=t_s,y=hw_filter]{figdata/networks/04-time-synchronisation/offset.csv};
\addplot[omDef,very thick] table[x=t_s,y=hw_tc]{figdata/networks/04-time-synchronisation/offset.csv};
\addplot[gray,dashed,domain=0:600] {100000};
\legend{horodatages logiciels,horodatages matériels,matériel et filtre de délai minimal,matériel et horloges transparentes}
\end{axis}
\end{tikzpicture}

\omcaption{Simulation~: l'horloge d'un serveur (20~ppm de décalage, avec dérive) disciplinée sur un chemin de trois commutateurs, avec une mise en file moyenne de
\qty{5}{\micro\second} par commutateur et par sens, huit échanges par seconde. La mise en file d'attente, et non les horodatages, domine jusqu'à ce que
les \omterm{def:nw:time-synchronisation:bc}{horloges transparentes} la suppriment~; la ligne en pointillés est la limite de 100~microsecondes. Valeurs de modèle~; données~: \texttt{fig\_clock.py} sur
\texttt{firm.clocksync} (graine fixée).}\label{fig:nw:time-synchronisation:offsets}
\end{omfigure}

Le tutoriel exécute les quatre configurations de la \cref{fig:nw:time-synchronisation:offsets} pendant une heure. Après les six premières minutes, le plus grand
écart atteint \qty{26.3}{\micro\second} avec des horodatages logiciels et \qty{25.9}{\micro\second} avec des horodatages matériels~: l'horodatage matériel n'apporte guère
lorsque les messages attendent pendant des microsecondes. En conservant, parmi chaque huit échanges, celui dont l'aller-retour est le plus court
(l'idée du filtre d'horloge de NTP~: l'échange le moins retardé est aussi le moins mis en file), on le ramène à \qty{11.1}{\micro\second}.
Les \omterm{def:nw:time-synchronisation:bc}{horloges transparentes} l'amènent à \qty{25}{\nano\second}, soit mille fois mieux. Avec des \omterm{def:nw:time-synchronisation:bc}{horloges transparentes}, une asymétrie fixe de
\qty{2}{\micro\second} laisse l'horloge décalée de \qty{1}{\micro\second}, exactement comme le dit la \cref{prop:nw:time-synchronisation:offset}, et aucun
filtrage ne la voit.

\omcode{code/firm/clocksync/firm_clocksync.py}{75}{91}{La simulation~: l'oscillateur dérive, le décalage croît, un échange le mesure,
et l'asservissement corrige la phase et la fréquence.}\label{lst:nw:time-synchronisation:loop}

\begin{omfigure}
\begin{tikzpicture}
\begin{axis}[width=11.5cm,height=5.4cm,xmode=log,ymode=log,xmin=0.1,xmax=700,ymin=1e-11,ymax=2e-7,
  xlabel={temps de moyennage $\tau$ (s, log)},ylabel={$\sigma_y(\tau)$ (log)},
  tick label style={font=\small},label style={font=\small},grid=major,grid style={gray!25},
  legend columns=2,legend style={font=\footnotesize,at={(0.5,-0.3)},anchor=north,draw=none,fill=none,
  /tikz/every even column/.append style={column sep=8pt}},table/col sep=comma]
\addplot[omThm,very thick,mark=*,mark size=1.4pt] table[x=tau_s,y=free]{figdata/networks/04-time-synchronisation/adev.csv};
\addplot[omDef,very thick,mark=square*,mark size=1.4pt] table[x=tau_s,y=disciplined]{figdata/networks/04-time-synchronisation/adev.csv};
\legend{oscillateur libre,horloge disciplinée (horloges transparentes)}
\end{axis}
\end{tikzpicture}

\omcaption{Simulation~: \omterm{def:nw:time-synchronisation:adev}{écart-type d'Allan} de l'oscillateur du chapitre laissé à lui-même et de la même horloge disciplinée sur des \omterm{def:nw:time-synchronisation:bc}{horloges
transparentes}. La fréquence de l'horloge libre dérive (pente $+\tfrac12$)~; les lectures de l'horloge disciplinée présentent une gigue de quelques nanosecondes autour d'une
cadence correcte (pente $-1$). En dessous de quelques secondes (entre 2 et 8), l'oscillateur libre est le plus stable des deux~; au-delà, l'asservissement l'emporte de plusieurs ordres de grandeur.
Valeurs de modèle~; données~: \texttt{fig\_clock.py}.}\label{fig:nw:time-synchronisation:adev}
\end{omfigure}

\section{Exigences réglementaires et preuve de l'horloge}

\begin{dated}{2026-09}{Ce que les règles exigent des horloges}\label{dat:nw:time-synchronisation:rules}
Dans l'Union européenne, la Commission Delegated Regulation (EU) 2017/574 («~RTS~25~») fixe, pour les membres et les participants des plateformes de négociation
qui recourent à une technique de trading algorithmique à haute fréquence, un écart maximal à l'UTC de 100 microsecondes et une granularité des horodatages de
1 microseconde ou mieux~; pour les plateformes dont la latence passerelle à passerelle est d'une milliseconde ou moins, la même exigence~; pour les systèmes vocaux et
les systèmes de demande de cotation (RFQ) à intervention humaine, une seconde. L'UTC provient d'un centre de temps inscrit à la liste du BIPM ou d'un système satellitaire
dont l'écart a été retiré. Les sociétés doivent établir la \omterm{def:nw:time-synchronisation:trace}{traçabilité à l'UTC}, identifier le point exact auquel chaque horodatage est appliqué, et
revoir la conformité au moins une fois par an. Aux États-Unis, la FINRA Rule 6820 oblige les membres du secteur déclarant au
consolidated audit trail à synchroniser leurs horloges d'activité à moins de 50 millisecondes de l'horloge atomique du NIST (une seconde pour les ordres manuels),
avant l'ouverture de chaque journée de négociation et tout au long de la journée, consigner la synchronisation pendant cinq ans, et certifier la conformité.
\end{dated}

\begin{definition}[Traçabilité à l'UTC]\label{def:nw:time-synchronisation:trace}
La \emph{traçabilité à l'UTC}\index{traçabilité à l'UTC} est une chaîne ininterrompue et documentée de comparaisons, allant des horodatages d'une société jusqu'à l'UTC,
chacune assortie d'une incertitude déclarée, de sorte que la société puisse montrer, pour tout horodatage, de combien il aurait pu s'écarter de l'UTC.
\end{definition}

\begin{method}[Prouver son horloge]\label{met:nw:time-synchronisation:prove}
\begin{enumerate}
\item Documenter la chaîne (\cref{fig:nw:time-synchronisation:chain})~: la référence, les \omterm{def:nw:time-synchronisation:gm}{horloges grand maître}, chaque \omterm{def:nw:time-synchronisation:bc}{horloge frontière} et transparente, l'horloge
  de la carte de chaque serveur, le passage à l'horloge de l'hôte et le point exact où chaque type d'horodatage est pris.
\item Établir le bilan d'erreur maillon par maillon, avec ses justificatifs (une spécification, une calibration, une mesure), et comparer son cas le plus
  défavorable à la limite.
\item Surveiller en continu~: le décalage mesuré et le délai de chemin de chaque esclave, l'état de verrouillage et le \omterm{def:nw:time-synchronisation:holdover}{maintien (holdover)} de l'\omterm{def:nw:time-synchronisation:gm}{horloge grand maître}, le passage de la carte à l'hôte~;
  déclencher une alarme bien en deçà de la limite.
\item Conserver les enregistrements pendant les durées de rétention exigées par les régulateurs, y compris chaque période en maintien ou hors tolérance.
\item Calibrer les asymétries fixes par rapport à une référence indépendante, et revérifier après chaque changement de câblage ou d'équipement.
\item Revoir l'ensemble du système au moins une fois par an.
\end{enumerate}
\end{method}

\section{Tutoriel~: discipliner une horloge}\label{tut:nw:time-synchronisation:tut}

\begin{tutorial}
\textbf{Objectif.} Discipliner de quatre façons une horloge de serveur simulée sur un réseau chargé, et lire ce qui limite chacune.
\textbf{État final~:} \cref{fig:nw:time-synchronisation:offsets,fig:nw:time-synchronisation:adev}, les quatre plus grands décalages cités dans le
  texte, et l'expérience sur l'asymétrie.
\begin{tutsteps}
\item \textbf{Les composants.} \texttt{firm.clocksync} regroupe l'oscillateur, le chemin (délai de base, asymétrie, file d'attente par commutateur, bruit d'horodatage),
  un échange (\cref{lst:nw:time-synchronisation:exchange}) et l'asservissement~; \texttt{simulate} exécute la boucle
  (\cref{lst:nw:time-synchronisation:loop}).
\item \textbf{Quatre configurations.} \texttt{nw\_clock.SETUPS} ne diffère que par le bruit d'horodatage, le filtrage et les \omterm{def:nw:time-synchronisation:bc}{horloges transparentes}~;
  \texttt{summary(name)} donne, pour chacune, le plus grand décalage après six minutes.
\item \textbf{Stabilité.} \texttt{adev()} calcule l'\omterm{def:nw:time-synchronisation:adev}{écart-type d'Allan} de l'horloge libre et de l'horloge disciplinée.
\item \textbf{Le point aveugle.} \texttt{asymmetry\_bias(2000)} ajoute une asymétrie fixe de \qty{2}{\micro\second} et renvoie le décalage
  qu'elle laisse, à comparer à $-a/2$. \texttt{python fig\_clock.py} écrit les données des graphiques.
\end{tutsteps}
\textbf{Pour aller plus loin.} Porter la file d'attente à \qty{50}{\micro\second} par commutateur et repérer les configurations qui
  dépassent la limite de 100 microsecondes~; diviser par deux le débit d'échange et relire le décalage de l'\omterm{def:nw:time-synchronisation:bc}{horloge transparente}.
\end{tutorial}

\section{Construction~: le modèle de synchronisation d'horloges}\label{bld:nw:time-synchronisation:clocksync}

\begin{build}
\bfield{Rôle} Un modèle de la chaîne de temporisation de la société, dans lequel l'effet de chaque choix (horodatages, filtres, \omterm{def:nw:time-synchronisation:bc}{horloges transparentes},
  asymétrie, maintien) peut être calculé avant l'achat, et un rapport de conformité des décalages par rapport à une règle.
\bfield{Interface} \texttt{firm\_clocksync}~: \texttt{Oscillator(freq\_ppm, walk\_ppb, seed)},
\texttt{Path(base\_ns, asym\_ns, queue\_ns, hops, stamp\_ns, seed)}, \texttt{Servo(kp, ki).update(offset, interval)}, \texttt{exchange},
\texttt{simulate(hours, interval\_s, osc, path, transparent, kp, ki, seed, min\_filter)}, \texttt{allan\_deviation}, \texttt{holdover\_s}, \texttt{compliance}~;
  jumeau C++20 de l'asservissement dans \texttt{cpp/}.
\bfield{Règles} Temps en nanosecondes, fréquences en ppb~; l'estimation est la formule à deux voies et rien d'autre~; les \omterm{def:nw:time-synchronisation:bc}{horloges transparentes}
  ne corrigent que la file d'attente~; chaque tirage est initialisé à partir d'une graine.
\bfield{Tests d'acceptation} \texttt{code/firm/clocksync/tests/}~: la formule à deux voies exacte sans asymétrie et biaisée de la moitié de celle-ci avec asymétrie~;
  les \omterm{def:nw:time-synchronisation:bc}{horloges transparentes} suppriment la file d'attente~; l'asservissement converge à partir d'une erreur de 20~ppm~; le montage de l'asservissement reproduit en Python et en C++~; un bruit de phase blanc
  donne une pente d'Allan de $-1$~; maintien avec et sans vieillissement~; comptage de conformité.
\bfield{Pour aller plus loin} Une \omterm{def:nw:time-synchronisation:bc}{horloge frontière} dotée de son propre asservissement entre le maître et l'esclave, et l'erreur qu'elle ajoute~; le mécanisme de délai entre pairs~; un
  rapport de surveillance dans le format qu'un régulateur demanderait à voir.
\end{build}

\begin{omsources}
\item Commission Delegated Regulation (EU) 2017/574 (RTS~25), articles 1 et 4 et annexe~; FINRA Rule 6820.
\item RFC~5905, \emph{Network Time Protocol Version 4}~; linuxptp, \texttt{ptp4l(8)}.
\item U.S. Department of Defense, \emph{GPS Standard Positioning Service Performance Standard}, 5th edition (2020).
\end{omsources}

\section{Exercices}

\begin{exercise}[$\star$]\label{exo:nw:time-synchronisation:1}
Une horloge avance de 15~ppm. Au bout de combien de temps après avoir été réglée correctement dépasse-t-elle une limite de 100 microsecondes, puis une limite de 50 millisecondes~?
\end{exercise}

\begin{exercise}[$\star$]\label{exo:nw:time-synchronisation:2}
Un échange donne $t_1 = 0$, $t_2 = 5\,300$, $t_3 = 6\,300$, $t_4 = 8\,900$ (nanosecondes). Calculer le décalage estimé
et le temps d'aller-retour.
\end{exercise}

\begin{exercise}[$\star$]\label{exo:nw:time-synchronisation:3}
La fibre allant du maître à l'esclave est plus longue de 40~mètres que la fibre du retour. Quel décalage la méthode à deux voies laisse-t-elle, avec un indice
de groupe de 1.462~?
\end{exercise}

\begin{exercise}[$\star\star$]\label{exo:nw:time-synchronisation:4}
Pourquoi les horodatages matériels ont-ils à peine aidé dans la \cref{fig:nw:time-synchronisation:offsets}, et pourquoi garder l'échange le moins retardé
sur huit aide-t-il~?
\end{exercise}

\begin{exercise}[$\star\star$]\label{exo:nw:time-synchronisation:5}
Une \omterm{def:nw:time-synchronisation:gm}{horloge grand maître} en maintien présente une erreur de fréquence résiduelle de 1~ppb et vieillit de $10^{-5}$~ppb par seconde. Combien de temps reste-t-elle
dans la limite de \qty{100}{\micro\second}~?
\end{exercise}

\begin{exercise}[$\star\star$]\label{exo:nw:time-synchronisation:6}
Lire la \cref{fig:nw:time-synchronisation:adev}~: pour quelles durées de moyennage l'horloge disciplinée est-elle moins bonne que l'horloge libre, et pourquoi
cela ne pose-t-il pas problème~?
\end{exercise}

\begin{exercise}[$\star\star\star$]\label{exo:nw:time-synchronisation:7}
\emph{Programmation.} Avec \texttt{firm.clocksync}, porter la file d'attente moyenne à \qty{50}{\micro\second} par commutateur et par sens, et calculer le
plus grand décalage après six minutes avec des horodatages matériels, avec le filtre à délai minimal et avec les \omterm{def:nw:time-synchronisation:bc}{horloges transparentes}. Lesquels
enfreignent la règle des 100 microsecondes~?
\end{exercise}

\begin{exercise}[$\star\star\star$]\label{exo:nw:time-synchronisation:8}
\emph{Trouver la faille.} «~Les statistiques PTP de nos serveurs montrent des décalages inférieurs à 50~ns toute la journée, donc nos horodatages sont à moins de 50~ns de l'UTC.~»
\end{exercise}

\section{Problème~: cent microsecondes jusqu'à l'UTC}

\begin{problem}[{Problème du week-end --- un bilan d'erreur et un maintien}]\label{pb:nw:time-synchronisation:1}
La chaîne de temporisation d'une société présente les termes d'erreur suivants (valeurs de modèle, en nanosecondes)~: le décalage UTC diffusé par le satellite 30~; le délai du câble d'antenne
après compensation 20~; l'\omterm{def:nw:time-synchronisation:gm}{horloge grand maître} 40~; deux \omterm{def:nw:time-synchronisation:bc}{horloges frontières} de 50 chacune~; l'asymétrie résiduelle du chemin après les \omterm{def:nw:time-synchronisation:bc}{horloges transparentes},
égale à la moitié de \qty{200}{\nano\second}~; l'horloge de la carte 20~; le passage de l'horloge de la carte à l'horloge de l'hôte 500~; le point d'horodatage
logiciel 1\,500. La règle fixe 100 microsecondes.

\medskip\noindent\textbf{Part I --- The budget.}
\begin{enumerate}
\item Quelle est l'erreur dans le pire cas, en sommant les termes~?
\item Quelle est la racine de la somme des carrés, si les termes sont indépendants~?
\item Quels sont les deux termes dominants, et quelle part du pire cas représentent-ils~?
\item Quelle marge le pire cas laisse-t-il par rapport à la règle~?
\end{enumerate}

\medskip\noindent\textbf{Part II --- Better timestamps.}
\begin{enumerate}[resume]
\item La société horodate les ordres sur la carte au lieu de le faire dans le logiciel. Quels sont le nouveau pire cas et la nouvelle racine de la somme des carrés~?
\item Pourquoi cela retire-t-il aussi le passage de la carte à l'hôte du bilan de ces horodatages~?
\item Quel terme domine désormais~?
\item Qu'est-ce qui diviserait par deux le pire cas restant~?
\end{enumerate}

\medskip\noindent\textbf{Part III --- Holdover.}
\begin{enumerate}[resume]
\item L'antenne tombe en panne. L'oscillateur de l'\omterm{def:nw:time-synchronisation:gm}{horloge grand maître} présente une erreur de fréquence résiduelle de 1~ppb et vieillit de $10^{-5}$~ppb par seconde. Combien de temps faut-il pour que la seule \omterm{def:nw:time-synchronisation:gm}{horloge grand maître} épuise la marge de la question~4~?
  Combien de temps faut-il pour que la seule \omterm{def:nw:time-synchronisation:gm}{horloge grand maître} épuise la marge de la question~4~?
\item Une \omterm{def:nw:time-synchronisation:gm}{horloge grand maître} meilleur marché présente une erreur résiduelle de 50~ppb et vieillit de $10^{-3}$~ppb par seconde. Combien de temps~?
\item Sur un week-end de 60~heures sans réparation, laquelle reste conforme~?
\item Que doit faire la supervision pendant le maintien~?
\end{enumerate}

\medskip\noindent\textbf{Part IV --- The verdict.}
\begin{enumerate}[resume]
\item Énoncer le \emph{résultat nommé}~: l'erreur dans le pire cas de la chaîne avec un horodatage logiciel et avec un horodatage sur carte, et les deux durées de maintien.
\item Comment l'exigence de granularité de la règle européenne contraint-elle l'horodatage logiciel~?
\item Pourquoi une seconde \omterm{def:nw:time-synchronisation:gm}{horloge grand maître}, avec une antenne indépendante, vaut-elle davantage qu'un meilleur oscillateur~?
\item Que exige la traçabilité au-delà du bilan~?
\item Lequel de ces termes un auditeur pourrait-il vérifier sans faire confiance à l'horloge propre de la société~?
\item Que ferait une asymétrie de 20~microsecondes (une liaison mal configurée) au bilan~?
\item Comment la règle américaine se compare-t-elle à ce bilan~?
\item En une phrase~: que doit savoir une société pour dire à quelle distance se trouvent ses horodatages de l'UTC~?
\end{enumerate}
\end{problem}

\section{Exercices}

\begin{interviewq}[$\star$ \iqroles{développeur}]\label{iq:nw:time-synchronisation:1}
Dérivez l'estimation de décalage à deux voies. Quelle hypothèse fait-elle~?
\end{interviewq}

\begin{interviewq}[$\star\star$ \iqroles{développeur}]\label{iq:nw:time-synchronisation:2}
Quelle est la différence entre une \omterm{def:nw:time-synchronisation:bc}{horloge frontière} et une \omterm{def:nw:time-synchronisation:bc}{horloge transparente}, et pourquoi chacune compte-t-elle dans un réseau de trading~?
\end{interviewq}

\begin{interviewq}[$\star\star$ \iqroles{développeur}]\label{iq:nw:time-synchronisation:3}
Pourquoi choisiriez-vous PTP plutôt que NTP pour un hôte de trading~? Quand NTP suffit-il~?
\end{interviewq}

\begin{interviewq}[$\star\star$ \iqroles{développeur, chercheur}]\label{iq:nw:time-synchronisation:4}
Vous comparez les horodatages de deux serveurs, et l'un indique qu'une exécution est arrivée avant l'ordre qui l'a causée. Que vérifiez-vous~?
\end{interviewq}

\begin{interviewq}[$\star\star\star$ \iqroles{développeur}]\label{iq:nw:time-synchronisation:5}
Comment prouveriez-vous à un régulateur que vos horodatages étaient à moins de 100~microsecondes de l'UTC mardi dernier~?
\end{interviewq}

\begin{interviewq}[$\star\star$ \iqroles{développeur}]\label{iq:nw:time-synchronisation:6}
L'antenne satellite tombe en panne un vendredi soir. Que devient votre horloge, et que faites-vous~?
\end{interviewq}
